%%%%%%%%%%%%%%%%%%%%%%%%%%%%%%%%%%%%%%%%%
% Medium Length Graduate Curriculum Vitae
% LaTeX Template
% Version 1.1 (9/12/12)
%
% This template has been downloaded from:
% http://www.LaTeXTemplates.com
%
% Original author:
% Rensselaer Polytechnic Institute (http://www.rpi.edu/dept/arc/training/latex/resumes/)
%
% Important note:
% This template requires the res.cls file to be in the same directory as the
% .tex file. The res.cls file provides the resume style used for structuring the
% document.
%
%%%%%%%%%%%%%%%%%%%%%%%%%%%%%%%%%%%%%%%%%

%----------------------------------------------------------------------------------------
%	PACKAGES AND OTHER DOCUMENT CONFIGURATIONS
%----------------------------------------------------------------------------------------

\documentclass[margin]{res} % The resume class defaults to 10pt.

\usepackage{helvet} % Default font is the helvetica postscript font
%\usepackage{newcent} % To change the default font to the new century schoolbook postscript font uncomment this line and comment the one above

\usepackage{amssymb}
\usepackage[T1]{fontenc}
\usepackage[utf8]{inputenc}



\usepackage{hyperref}
\hypersetup{
    colorlinks=true,
    linkcolor=blue,
    filecolor=magenta,      
    urlcolor=blue,
}
\urlstyle{same}

\usepackage{etaremune}

\setlength{\textwidth}{5.1in} % Text width of the document
\setlength{\emergencystretch}{1em} % Allow long publication titles to wrap cleanly.
\interlinepenalty=10000 % Keep each publication or talk paragraph on one page.


\begin{document}

%----------------------------------------------------------------------------------------
%	NAME AND ADDRESS SECTION
%----------------------------------------------------------------------------------------

\moveleft.5\hoffset\centerline{\large\bf Mahesh Sreekumar Rajasree\footnote{\href{mailto:maheshsr@mimuw.edu.pl}{maheshsr@mimuw.edu.pl}}\textsuperscript{\color{blue},}\footnote{\href{https://mahe94.github.io}{https://mahe94.github.io}}} % Your name at the top
 
\moveleft\hoffset\vbox{\hrule width\resumewidth height 1pt}\smallskip % Horizontal line after name; adjust line thickness by changing the '1pt'
 
\moveleft.5\hoffset\centerline{Assistant Professor, University of Warsaw}
\moveleft.5\hoffset\centerline{Institute of Informatics}
\moveleft.5\hoffset\centerline{Faculty of Mathematics, Informatics and Mechanics}
\moveleft.5\hoffset\centerline{Warsaw, Poland}
%\moveleft.5\hoffset\centerline{+91-96455452535}

%----------------------------------------------------------------------------------------

\begin{resume}

 
\section{PRINCIPAL \\ INTERESTS}  

Cryptography, Quantum Cryptography, Lattice theory, Algorithms, Cryptanalysis


\section{Employment}

{\sl Assistant Professor} \hfill October 2026 - Present \\
Institute of Informatics\\
\href{https://www.mimuw.edu.pl/en/}{Faculty of Mathematics, Informatics and Mechanics}\\
University of Warsaw, Warsaw, Poland

{\sl Postdoctoral Researcher} \hfill December 2024 - September 2026 \\
\textbf{Host:} \href{https://sites.google.com/site/nicodoettling/}{Prof. Nico Döttling} \\
\href{https://cispa.de/en}{CISPA Helmholtz Center for Information Security}\\
Saarbrücken, Saarland, Germany


{\sl Postdoctoral Fellow} \hfill April 2023 - November 2024 \\
\textbf{Supervisor:} \href{https://web.iitd.ac.in/~kvenkata/}{Prof. Venkata Koppula}\\
\href{https://homecse.iitd.ac.in}{Computer Science and Engineering}\\
Indian Institute of Technology Delhi, New Delhi, India


 
\section{ACADEMIC \\BACKGROUND}

{\sl Ph.D. \& M.Tech. \href{https://www.cse.iitk.ac.in}{Computer Science and Engineering}} \hfill August 2016 - July 2023 \\
Indian Institute of Technology Kanpur, Uttar Pradesh, India
\begin{itemize}
\item Advisor: \href{https://en.wikipedia.org/wiki/Manindra_Agrawal}{Prof. Manindra Agrawal}
\item Ph.D. Thesis: \href{https://mahe94.github.io/assets/bibliography/phd_thesis.pdf}{Analysis of Symmetric-Key Cryptosystems and Subset-Sum Problem}
\item M.Tech. Thesis: \href{https://mahe94.github.io/assets/bibliography/mtech_thesis.pdf}{New Results and Proofs in Lattice Theory}
\item CPI: 10/10
\end{itemize} 

{\sl B.Tech. \href{https://nitc.ac.in/department/computer-science-amp-engineering}{Computer Science and Engineering}} \hfill August 2012 - May 2016 \\
National Institute of Technology Calicut, Kerala, India
\begin{itemize}
\item Advisor: \href{https://nitc.ac.in/department/computer-science-amp-engineering/faculty-and-staff/faculty/a8fddc9d-779e-4109-9846-3e0198986145}{Dr. Sumesh T.A.}
\item B.Tech. Thesis: \href{https://mahe94.github.io/assets/bibliography/btech_thesis.pdf}{Firewall for Preventing Anonymous Proxy Usage}
\item CPI: 8.89/10
\end{itemize} 

{\sl Secondary \& Higher Secondary Education} \hfill June 2009 - March 2012 \\
\href{https://www.aryacentralschool.org}{Arya Central School}, Trivandrum, Kerala, India
\begin{itemize}
    \item All India Senior Secondary School Examination \qquad Percentage:  91.2
    \item All India Secondary School Examination  \qquad CGPA: 9.6/10
\end{itemize}



\section{Peer-reviewed Works}
\begin{enumerate}

\item \textbf{On Robust Combiners for Pseudorandom Codes and PKE with Pseudorandom Ciphertexts}\\
with Nico Döttling, Venkata Koppula and Felix Rohrbach\\
\emph{Theory of Cryptography Conference, TCC 2026}

\item \textbf{Pigeonhole Equal Subset Sum: Subexponential Algorithm, Tight Lower Bounds and Average-Case Analysis}\\
with Deepak Bhati, Pranjal Dutta, Antoine Joux and Karol W\k{e}grzycki\\
\emph{67th Annual Symposium on Foundations of Computer Science, FOCS 2026}

\item \textbf{Chosen Ciphertext Secure Pseudorandom Codes in the Standard Model}\\
with Nico Döttling, Antoine Joux, Venkata Koppula and Hendrik Waldner\\
\emph{46th Annual International Cryptology Conference, CRYPTO 2026}

\item \textbf{Separating Pseudorandom Codes from Local Oracles}\\
with Nico Döttling and Anne Müller\\
\emph{23rd Theory of Cryptography, TCC 2025}\\
\textbf{\textcolor{blue}{Outstanding Paper Award}}

\item \textbf{Incompressible Encryption Beyond CPA/CCA Security}\\
with Venkata Koppula, Abhinav Kumar and Harihar Swaminathan\\
\emph{IACR Communications in Cryptology, Volume 2, Issue 2, Article 5, 2025}

\item \textbf{A Note on Adaptive Security in Hierarchical Identity-Based Encryption}\\
with Rishab Goyal and Venkata Koppula\\
\emph{45th Annual International Cryptology Conference, CRYPTO 2025}

\item \textbf{Non-Committing Identity Based Encryption: Constructions and Applications}\\
with Rishab Goyal, Fuyuki Kitagawa, Venkata Koppula, Ryo Nishimaki and Takashi Yamakawa\\
\emph{28th IACR International Conference on Practice and Theory of Public-Key Cryptography, PKC 2025}
    
\item \textbf{Incompressible Functional Encryption}\\
with Rishab Goyal, Venkata Koppula and Aman Verma\\
\emph{16\textsuperscript{th} Innovations in Theoretical Computer Science, ITCS 2025}

\item \textbf{Leakage-Resilient Key-Dependent Message\\Secure Encryption Schemes}\\
with Dhairya Gupta and Harihar Swaminathan\\
\emph{25\textsuperscript{th} International Conference on Cryptology, INDOCRYPT 2024}    

\item \textbf{Leakage-Resilient Incompressible Cryptography: Constructions and Barriers}\\
with Kaartik Bhushan, Rishab Goyal, Venkata Koppula, Varun Narayanan and Manoj Prabhakaran,\\
\emph{International Conference on the Theory and Application of Cryptology and Information Security, ASIACRYPT 2024}

\item \textbf{Complexity of Monomial Prediction in Cryptography and Machine Learning}\\
with Pranjal Dutta and Santanu Sarkar,\\
\emph{26\textsuperscript{th} International Symposium on Symbolic and Numeric Algorithms for Scientific Computing, SYNASC 2024}

\item \textbf{On the Maximum Distance Sublattice Problem and Closest Vector Problem}\\
with Rajendra Kumar and Shashank K Mehta,\\
\emph{26\textsuperscript{th} International Symposium on Symbolic and Numeric Algorithms for Scientific Computing, SYNASC 2024}

\item \textbf{Efficient reductions and algorithms
for Subset Product}\\
with Pranjal Dutta,\\
\emph{9\textsuperscript{th}  Annual International Conference on Algorithms and Discrete Applied Mathematics, CALDAM 2023}\\
\textbf{\textcolor{blue}{Best Student Presentation Award}}

\item \textbf{On the Variants of Subset Sum: Projected and Unbounded}\\
with Pranjal Dutta,\\
\emph{25\textsuperscript{th} International Symposium on Symbolic and Numeric Algorithms for Scientific Computing, SYNASC 2023}

\item \textbf{On the bases of $\mathbb{Z}^n$ lattice}\\
with Shashank K Mehta,\\
\emph{24\textsuperscript{th} International Symposium on Symbolic and Numeric Algorithms for Scientific Computing, SYNASC 2022}

\item \textbf{Weak-keys and key-recovery attack for TinyJAMBU}\\
with Pranjal Dutta and Santanu Sarkar,\\
\emph{Scientific Reports, Volume 12, Article 16313, 2022}

\item \textbf{On the hardness of monomial prediction and zero-sum distinguishers for Ascon}\\
with Pranjal Dutta and Santanu Sarkar,\\
\emph{12\textsuperscript{th} International Workshop on Coding and Cryptography, WCC 2022}

\item \textbf{Algebraic algorithms for variants of Subset Sum}\\
with Pranjal Dutta,\\
\emph{8\textsuperscript{th}  Annual International Conference on Algorithms and Discrete Applied Mathematics, CALDAM 2022}\\
\textbf{\textcolor{blue}{Best Student Presentation Award}}

\item \textbf{Cryptanalysis of Round-Reduced KECCAK using Non-Linear Structures}\\
\emph{20\textsuperscript{th} International Conference on Cryptology, INDOCRYPT 2019}

\item \textbf{Cryptanalysis of 1-Round KECCAK}\\
with Rajendra Kumar and Hoda Al Khzaimi,\\
\emph{10\textsuperscript{th} International Conference on Cryptology, AFRICACRYPT 2018}

\end{enumerate}



\section{Papers Under Submission}

\begin{enumerate}
\item \textbf{On the Largest Cardinality Subset Sum}\\
with Deepak Bhati, Anubhav Dhar, Pranjal Dutta and Karol W\k{e}grzycki

\item \textbf{A Simple Compiler for CCA2-Secure Pseudorandom Codes in the Standard Model}\\
with Nico Döttling, Antoine Joux, Venkata Koppula and Hendrik Waldner

\item \textbf{Weak Pigeonhole Equal-Sums made Simpler and Faster}\\
with Deepak Bhati, Pranjal Dutta, Antoine Joux and Karol W\k{e}grzycki

\item \textbf{Quantum Pseudorandom Codes}\\
with Nico Döttling and Anne Müller

\end{enumerate}




% \section{Unpublished works (Not submitted)}

% \begin{enumerate}
    
% \end{enumerate}


\section{Academic Achievements}

\begin{enumerate}
    \item Won the \href{https://www.linkedin.com/posts/cispa_congratulations-on-the-tcc-outstanding-ugcPost-7407070985310580737-7lt2?utm_source=share&utm_medium=member_desktop&rcm=ACoAACTz6fQBs27BwJxjG_fen9iqLzLlrNMFprs}{Outstanding Paper award} in TCC 2025.
    
    \item Won the \href{http://acsdm.in/acsdm/awards/}{Best Student Presentation award} (one of the 5 winners) in CALDAM 2023.
    
    \item Won the \href{http://acsdm.in/acsdm/awards/}{Best Student Presentation award} (one of the 3 winners) in CALDAM 2022.

    \item Awarded \href{https://pmrf.in}{\textbf{Prime Minister’s Research Fellowship}} in 2018 (received approximately USD 50,000 for a duration of 4 years to cover stipend and other expenses related to research).
    
    \item Awarded \href{https://mahe94.github.io/assets/pdf/excellence_certificate.pdf}{Certificate of Merit for Academic Excellence} in M.Tech. Computer Science and Engineering at IIT Kanpur, academic year 2016--2017.
    
    \item Secured \href{https://mahe94.github.io/assets/pdf/jest.pdf}{Rank 28 (Computer Science, Part A)} in \href{https://en.wikipedia.org/wiki/Joint_Entrance_Screening_Test}{Joint Entrance Screening Test} (JEST-2016).
    
    \item Secured \href{https://mahe94.github.io/assets/pdf/gate.pdf}{All India Rank 575} in \href{https://en.wikipedia.org/wiki/Graduate_Aptitude_Test_in_Engineering}{Graduate Aptitude Test in Engineering} (GATE 2016)
    
    \item Secured fourth place in the first round of \href{https://inctf.in}{Amrita InCTF} 2015.
\end{enumerate}



\section{Short Term Research Visit}

\begin{enumerate}
    \item \textbf{National Institute of Technology Calicut} (Aug 2026), Host: Prof. K. Murali Krishnan.
    \item \textbf{Indian Institute of Space Science and Technology, Trivandrum} (Aug 2026), Host: Prof. Deepak Mishra.
    \item \textbf{Universidad Carlos III de Madrid, Leganés} (Jun 2026), Host: Prof. Lorena González-Manzano.
    \item \textbf{University of Warsaw} (Mar 2026).
    \item \textbf{IMDEA Software Institute} (May 2026), Host: \href{https://www.dariofiore.it}{Prof. Dario Fiore}
    
    \item \textbf{Technical University of Denmark} (Dec 2025), Host: \href{https://www.christianmajenz.info/about-me.html}{Prof. Christian Majenz}
    
    \item \textbf{International Institute of Information Technology Hyderabad} (Aug 2024), Host: \href{https://suryajith.github.io}{Dr. Suryajith Chillara}

    \item \textbf{University of Wollongong, Australia} (Apr 2024), AICCTP workshop visit.
    \item \textbf{Indian Institute of Space Science and Technology, Trivandrum} (Mar 2024), Host: Prof. Deepak Mishra.
    \item \textbf{Microsoft Research India, Bangalore} (Oct 2023), Host: Satya Lokam.
    \item \textbf{International Centre for Theoretical Sciences Bangalore} (Oct 2023), Host: \href{https://en.wikipedia.org/wiki/Jaikumar_Radhakrishnan}{Prof. Jaikumar Radhakrishnan}

    \item \textbf{Tata Institute of Fundamental Research, Mumbai} (Jul 2023), Host: Prof. Vinod M. Prabhakaran.
    \item \textbf{Indian Institute of Technology Bombay} (Jul 2023), Host: \href{https://www.cse.iitb.ac.in/~mp/}{Prof. Manoj M. Prabhakaran}

    \item \textbf{Institute of Science and Technology Austria} (Sept 2022), Host: \href{https://scholar.google.com/citations?user=uEjyGZUAAAAJ&hl=en}{Prof. Krzysztof Pietrzak}
\end{enumerate}

\section{Research Internships \& Projects}

\begin{enumerate}
    \item Visiting scholar at IIT Delhi hosted by \href{https://web.iitd.ac.in/~kvenkata/}{Prof. Venkata Koppula} (Feb 2023).

    \item Visiting scholar at Microsoft Research India hosted by \href{https://www.microsoft.com/en-us/research/people/satya/}{Satya Lokam} (Jun 2018, Oct 2018, Mar 2019, Sept 2019).
    
    \item Visiting scholar at New York University, Abu Dhabi hosted by \href{https://nyuad.nyu.edu/en/academics/divisions/engineering/faculty/hoda-al-khzaimi.html}{Prof. Hoda AlKhzaimi} (Oct 2017).
    
%     \item Automated subglottic secretions drainage device for patients in ventilators - Programmed the arduino of a device which helps to drain subglottic secretions automatically at particular intervals on
% the patient specific data input by the doctor.
\end{enumerate}


\section{Academic Talks and Presentations}

\begin{enumerate}
    \item \textbf{Recent Trends in Cryptography}, FSTTCS 2026 workshop, IIT Delhi, 19--20 December 2026. Scheduled speaker (upcoming).

    \item \textbf{Invited Talk}, National Institute of Technology Calicut, 31 August 2026.\\
    Title: Limits of Black-Box Constructions of Pseudorandom Codes.

    \item \textbf{Invited Talk}, IIST, Trivandrum, 18 August 2026.\\
    Title: Secrets, Quantum and AI.

    \item \textbf{Invited Talk}, IMDEA Software Institute, Madrid, 25 May 2026.\\
    Title: Limits of Black-Box Constructions of Pseudorandom Codes.

    \item \textbf{ITCS Graduating Bits}, Informal Talk.\\

    \item \textbf{CRYPTO 2025}, Conference Talk.\\
    Title: A Note on Adaptive Security in Hierarchical Identity-Based Encryption
    
    \item \textbf{PKC 2025}, Conference Talk.\\
    Title: Non-Committing Identity Based Encryption: Constructions and Applications
    
    \item \textbf{ITCS 2025}, Conference Talk.\\
    Title: Incompressible Functional Encryption

    \item \textbf{INDOCRYPT 2024}, Conference Talk.\\
    Title: Leakage-Resilient Key-Dependent Message Secure Encryption Schemes
    
    \item \textbf{SYNASC 2024}, Conference Talk.\\
    Title: Complexity of Monomial Prediction in Cryptography and Machine Learning

    \item \textbf{SYNASC 2024}, Conference Talk.\\
    Title: On the Maximum Distance Sublattice Problem and Closest Vector Problem
    
    \item \textbf{Invited Talk}, IIT Delhi, 2024.\\
    Title: Leakage-Resilience meets Incompressibility.

    \item \textbf{Continuing Education Programme}, Centre for Railway Information Systems employees, IIT Delhi, 2024. Guest lecture.

    \item \textbf{Invited Talk}, CSTAR, IIIT Hyderabad, 7 August 2024.\\
    Title: Leakage-Resilience meets Incompressibility
    
    \item \textbf{Invited Talk}, University of Wollongong, Australia, 3 April 2024.\\
    Title: Incompressible Encryption beyond Public Key Encryption
    
    \item \textbf{Invited Talk}, IIST, Trivandrum, 14 March 2024.\\
    Title: Incompressible Encryption beyond Public Key Encryption
    
    \item \textbf{Invited Talk}, Microsoft Research, Bangalore, 27 October 2023.\\
    Title: Incompressible Encryption beyond Public Key Encryption
    
    \item \textbf{SYNASC 2023}, Conference Talk.\\
    Title: On the Variants of Subset Sum: Projected and Unbounded
    
    \item \textbf{Invited Talk}, Trust Lab, IIT Bombay, 12 July 2023 (online).\\
    Title: Incompressible Encryption
    
    \item \textbf{CALDAM 2023}, Conference Talk.\\
    Title: Efficient reductions and algorithms for Subset Product

    \item \textbf{SYNASC 2022}, Conference Talk.\\
    Title: On the bases of $\mathbb{Z}^n$ lattice
    
    \item \textbf{SIGTACS}, CSE, IIT Kanpur.\\
    Title: On the hardness of monomial prediction problem.
    
    \item \textbf{STEM Mentorship Program for High School Girl Students}, IIT Delhi, 2023.\\
    Title: Cryptography: Performing a Task Securely.

    \item \textbf{Invited Talk}, CSE, IUST Jammu \& Kashmir, 2022.\\
    Title: On the Subset Sum Problem.
    
    \item \textbf{SIGTACS}, CSE, IIT Kanpur.\\
    Title: An Algebraic Algorithm for Hamming Subset-Sum Problem.
    
    \item \textbf{WCC 2022}, Workshop Talk.\\
    Title: On the hardness of monomial prediction problem.
    
    \item \textbf{Invited Talk}, PMRF Symposium, 2022.\\
    Title: Cryptanalysis of KECCAK \& Algorithms for Lattice Problems.
    
    \item \textbf{SIGTACS}, CSE, IIT Kanpur.\\
    Title: Cryptanalysis of Round-Reduced KECCAK using Non-linear Structures.
    
    \item \textbf{INDOCRYPT 2019}, Conference Talk.\\
    Title: Cryptanalysis of Round-Reduced KECCAK using Non-Linear Structures.
    
    \item \textbf{SIGTACS}, CSE, IIT Kanpur.\\
    Title: Leftover Hash Lemma
    
    \item \textbf{SIGTACS}, CSE, IIT Kanpur.\\
    Title: Lattices in Computer Science
    
    \item \textbf{SIGTACS}, CSE, IIT Kanpur.\\
    Title: A Deterministic Exponential Time Algorithm for some lattice problems
    
    
\end{enumerate}

\section{Teaching Experience}
\begin{enumerate}
    \item \href{https://informatorects.uw.edu.pl/en/courses/view?prz_kod=1000-2M26QC}{1000-2M26QC -- Quantum Computing}, University of Warsaw, winter semester 2026--2027. Course coordinator and instructor for lectures, tutorials and laboratory sessions; taught in English (9 ECTS).

    \item \href{https://informatorects.uw.edu.pl/en/courses/view?prz_kod=1000-2M12KI1}{1000-2M12KI1 -- Cryptography I}, University of Warsaw, winter semester 2026--2027. Tutorial instructor for two class groups.

    \item \href{https://informatorects.uw.edu.pl/en/courses/view?prz_kod=1000-2D21KBF}{1000-2D21KBF -- Cryptography, Blockchain \& Fintech}, University of Warsaw, 2026--2027. Master's seminar co-coordinated with Prof. Stefan Dziembowski and Prof. Tomasz Michalak; seminar instructor.

    \item \href{https://cms.cispa.saarland/crypto25/}{CRYPTO25 - Cryptography} at Saarland University, 2025. Delivered a week-long lecture series on key agreement protocols and public-key encryption.
    
    \item \href{https://kumarmadhukar.github.io/courses/col100-summer24/index.html}{COL100 - Introduction to Computer Science} at IITD, 2024. Delivered lectures, designed examinations, and supervised laboratory sessions, covering approximately one-third of the course.
    
    \item \href{https://sites.google.com/view/nikhilbalaji/holi2024}{COL863 - Algebra and Computation (Topics in TCS)} at CSE, IITD, 2024. Delivered a module on lattices, designed assignments, and graded coursework.
\end{enumerate}

\section{Research Mentoring}
\begin{enumerate}
    \item PhD research mentoring:
        \begin{enumerate}
            \item \href{https://cispa.de/en/people/c01debh}{Deepak Bhati}, CISPA

            \item \href{https://cispa.de/en/people/c01moba}{Hadi Barzegar}, CISPA

            \item \href{https://anne-me.github.io}{Anne Müller}, CISPA
        \end{enumerate}
    
    \item Master and dual-degree research mentoring: 
        \begin{enumerate}
            \item \href{https://shankhgupta.github.io}{Shankh Gupta}, IIT Delhi, dual-degree research, 2024--2025

            \item \href{https://mahe94.github.io/assets/pdf/abhinav.pdf}{Abhinav Kumar}, IIT Delhi, 2023--2024 (co-authored paper in Communications in Cryptology, 2025)

            \item \href{https://mahe94.github.io/assets/pdf/dhairya.pdf}{Dhairya Gupta}, IIT Delhi, 2023--2024 (co-authored paper at INDOCRYPT 2024) 
        \end{enumerate}

    \item Bachelor project co-mentoring:
        \begin{enumerate}
            \item \href{https://mahe94.github.io/assets/pdf/naman.pdf}{Naman Verma}, IIT Kanpur, 2018--2019 (with Prof. Shashank K. Mehta).
        \end{enumerate}
    
    \item Bachelor projects mentored: 
        \begin{enumerate}
            \item \href{https://ab271202.github.io/notes}{Anish Banerjee}, IIT Delhi, 2024--2025


            \item \href{https://scholar.google.com/citations?user=WKfskVwAAAAJ&hl=en}{Harihar Swaminathan}, IIT Delhi, 2023--2024 (co-authored paper in Communications in Cryptology, 2025 and INDOCRYPT 2024) 
        \end{enumerate}
    \item Research interns mentored at IIST, 2024:
        \begin{enumerate}
            \item \href{https://mahe94.github.io/assets/pdf/sundar.pdf}{Sundar R}, B.Tech. student, IIIT Sri City (NTRU cryptography).
            \item \href{https://mahe94.github.io/assets/pdf/aiswarya.pdf}{Aiswarya G Raj}, M.Tech. student, CET Trivandrum (adversarial robustness of deep neural networks; co-mentored with Prof. Deepak Mishra).
            \item \href{https://mahe94.github.io/assets/pdf/ajay.pdf}{Ajay R Nair}, M.Tech. student, CET Trivandrum (natural adversarial examples; co-mentored with Prof. Deepak Mishra).
        \end{enumerate}

\end{enumerate}

% \section{Undergraduate \& Graduate Advising}
% \begin{enumerate}
%     \item Dhairya Gupta (Dual), IITD, 2023
%     \item Abhinav Kumar (Dual), IITD, 2023
%     \item Harihar Swaminathan (BTech), IITD, 2023
%     \item Aman Verma (Dual), IITD, 2023
%     \item Naman Verma (BTech), IITK, 2019
% \end{enumerate}

\section{Teaching Assistance}
\begin{enumerate}
    \item COL872 - Special Topics in Cryptography at CSE, IITD, 2023
    \item CS961 - Modern Cryptology at Emaster Program, IITK 2023, 2022
    \item CS972 - Introduction to Linear Algebra at Emaster Program, IITK 2022
    \item Modern Algebra at NPTEL 2021, 2018
    \item CS641 - Modern Cryptology at CSE, IITK 2021, 2020, 2019, 2018, 2017
    \item CS345 - Algorithms II at CSE, IITK 2018
    \item ESO207 - Data Structures and Algorithms at CSE, IITK 2017
    \item CS3092D - Operating Systems Laboratory at CSE, NITC 2016
\end{enumerate}


\section{Professional Activities}
\begin{enumerate}
    \item Co-organizer, \href{https://mahe94.github.io/workshops/FSTTCS-2026/}{Recent Trends in Cryptography}, FSTTCS 2026 workshop, IIT Delhi, with Prof. Venkata Koppula. Scheduled for 19--20 December 2026 (upcoming).

    \item Reviewed papers for INDOCRYPT 2021, MFCS 2023, FSTTCS 2023, Journal of Cryptology, EUROCRYPT 2023, ASIACRYPT 2024, ITCS 2025, EUROCRYPT 2025, PKC 2025, ICALP 2025, CRYPTO 2025, STOC 2026, ICALP 2026, TCC 2026, ASIACRYPT 2026.

    \item PhD Defence Committee Member, Saarland University, 19 March 2026. Invited examiner for the doctoral defense of \href{https://publikationen.sulb.uni-saarland.de/bitstream/20.500.11880/41623/1/Enhancing%20the%20Security%20and%20Efficiency%20of%20Distributed%20Key%20Generation%20and%20Its%20Applications.pdf}{Renas Bacho}, evaluated thesis quality and oral examination performance.

    \item Part of the Indian team for the \href{https://iitk.ac.in/institute-and-uow-cloud-computing}{Australia--India Cyber and Critical Technology Partnership} (AICCTP). The IIT Kanpur--University of Wollongong collaboration studies privacy-enhancing cryptography and cloud-computing standardization.
\end{enumerate}

\section{Volunteer Works}
\begin{enumerate}
    \item Lecture delivered at Continuing Education Programme for employees at Centre For Railway Information Systems, IIT Delhi 2024.
    \item Lecture delivered at \href{https://home.iitd.ac.in/show.php?id=397&in_sections=News}{STEM Mentorship Program for High School Girl Students}, IIT Delhi 2023.

    \item Volunteer at \href{https://www.fsttcs.org.in/archives/2017/#oc}{Foundations of Software Technology and Theoretical Computer Science}, FSTTCS 2017.
    
    \item Volunteer at \href{https://cse.iitk.ac.in/doc/CSAW'17%20agenda-1%20(002).pdf}{Cybersecurity Awareness Week} CSAW 2017.
			
	\item Head Coordinator of \href{https://sigtacs.github.io}{Special Interest Group on Theoretical Aspects of Computer Science} SIGTACS at IIT Kanpur.
\end{enumerate}


\section{Conference \& Workshop Attended (offline)}
\begin{enumerate}
    \item CISPA--MPI--UdS retreat, Trier, Germany, November 2025.
    \item 23\textsuperscript{rd} International Conference on Theory of Cryptography, TCC 2025, Aarhus, Denmark.

    \item 45\textsuperscript{th} Annual International Cryptology Conference, CRYPTO 2025, Santa Barbara, CA, USA.
    
    \item Cryptography 10 Years Later: Obfuscation, Proof Systems, and Secure Computation at Simons Institute for the Theory of Computing (2025), Berkeley, USA.
    
    \item 28\textsuperscript{th} IACR International Conference on Practice and Theory of Public-Key Cryptography, PKC 2025, Røros, Norway.


    \item 16\textsuperscript{th} Innovations in Theoretical Computer Science Conference, ITCS 2025, Columbia University, New York, USA.

    \item 22\textsuperscript{nd} International Conference on Theory of Cryptography, TCC 2024, Milan, Italy.
    
    \item Warsaw IACR Summer School on Post-Quantum Cryptography 2024, Warsaw, Poland. 
    
    \item 9\textsuperscript{th} Annual International Conference on Algorithms and Discrete Applied Mathematics, CALDAM 2023, Gandhinagar, India.
    
    \item 24\textsuperscript{th} International Symposium on Symbolic and Numeric Algorithms for Scientific Computing, SYNASC 2022, Hagenberg, Austria.
    
    \item Quantum Algorithms at Simons Institute for the Theory of Computing (2020), Berkeley, USA.
    
    \item Quantum Cryptanalysis of Post-Quantum Cryptography at Simons Institute for the Theory of Computing (2020), Berkeley, USA.
    
    \item Lattices: Geometry, Algorithms and Hardness at Simons Institute for the Theory of Computing (2020), Berkeley, USA.
    
    \item Secure Multiparty Computation: Theory and Practice workshop held at IISc (2020), Bangalore, India.
    
    \item 20\textsuperscript{th} International Conference on Cryptology, INDOCRYPT 2019, Hyderabad, India.
    
    \item Theory of Cryptography Conference, TCC 2018, Goa, India.
    
    \item 37\textsuperscript{th} Foundations of Software Technology and Theoretical Computer Science 2017, Kanpur, India.

\end{enumerate}

\section{Conference \& Workshop Attended (online)}
\begin{enumerate}
    \item 26\textsuperscript{th} International Symposium on Symbolic and Numeric Algorithms for Scientific Computing, SYNASC 2024, Timișoara, Romania.
    
    \item 25\textsuperscript{th} International Symposium on Symbolic and Numeric Algorithms for Scientific Computing, SYNASC 2023, Nancy, France.

    \item 12\textsuperscript{th} International Workshop on Coding and Cryptography, WCC 2022.
    
    \item 8\textsuperscript{th} Annual International Conference on Algorithms and Discrete Applied Mathematics, CALDAM 2022
\end{enumerate}

\section{Referees}
\begin{enumerate}
    \item \textbf{Prof. Nico Döttling}\\
    CISPA Helmholtz\\
    Email: doettling@cispa.de

    \item \textbf{Prof. Antoine Joux}\\
    CISPA Helmholtz\\
    Email: joux@cispa.de
    
    \item \textbf{Prof. Venkata Koppula}\\
    Department of CSE, IIT Delhi\\
    Email: kvenkata@cse.iitd.ac.in
    % Mob: +919380937954

    \item \textbf{Prof. Manoj Prabhakaran}\\
    Department of CSE, IIT Bombay\\
    Email: mp@cse.iitb.ac.in
    % Mob: +91-22-2576-7709

    \item \textbf{Prof. Manindra Agrawal}\\
    Department of CSE, IIT Kanpur\\
    Email: manindra@cse.iitk.ac.in
    % Mob: +919935062605

    % \item \textbf{Dr. Rishab Goyal}\\
    % Assistant Professor\\
    % Department of CSE, UW-Madison\\
    % Email: rishab.goyal@wisc.edu

    % \item \textbf{Dr. Nikhil Balaji}\\
    % Assistant Professor\\
    % Department of CSE, IIT Delhi\\
    % Email: nbalaji@cse.iitd.ac.in\\
    % Mob: +919962245830

    \item \textbf{Prof. Shashank K Mehta}\\
    Department of CSE, IIT Kanpur (Retired)\\
    Email: skmehta@cse.iitk.ac.in
    % Mob: +919044407835
\end{enumerate}

\end{resume}
\end{document}
